\section{Self-Consistency：从采样到投票}

\subsection{动机：LLM 的概率本质}

Denny Zhou 从第一性原理出发推导 Self-Consistency 的合理性：

\begin{itemize}
    \item LLM 在解码时优化的是：$\arg\max P(\text{reasoning path}, \text{answer} \mid \text{problem})$
    \item 但我们真正需要的是：$\arg\max P(\text{answer} \mid \text{problem})$
    \item 两者的关系：$P(\text{answer} \mid \text{problem}) = \sum_{\text{path}} P(\text{answer}, \text{path} \mid \text{problem})$
    \item 即需要对所有可能的推理路径求和（marginalize），而非只取单一最优路径。
\end{itemize}

\subsection{方法与效果}

\begin{importantbox}{Self-Consistency 方法}
\begin{enumerate}
    \item 对同一问题采样多次（使用非零温度），获得多条推理路径和对应答案。
    \item 取\textbf{出现频率最高的答案}作为最终预测（多数投票 / majority voting）。
    \item 注意：投票对象是\textbf{最终答案}，而非推理路径——推理路径是隐变量（latent variable）。
\end{enumerate}
\end{importantbox}

Self-Consistency 带来了巨大的性能提升，在当时大幅刷新了所有基准的 SOTA。更重要的是，当一致性超过 80\% 时，准确率接近 100\%——这提供了一个可靠的置信度指标。

\begin{knowledgebox}{与 Universal Self-Consistency 的扩展}
对于自由形式（free-form）答案，可以通过 LLM 自身来判断哪些答案语义等价，找出最常见的响应聚类。例如对于 ``哪些国家人均咖啡消费低于墨西哥'' 这类问题，不同回答可能措辞不同但内容一致。
\end{knowledgebox}

\subsection{本章小结}

Self-Consistency 将 LLM 推理从单次生成提升到统计推断层面。其核心思想源自概率图模型中的边际化推理（marginal inference），简单但极为有效。

